\section{Conclusão}

Mineramos as quatro métricas DORA em 300 repositórios populares do GitHub que usam GitHub Actions, em um ano, com duas variantes do lead time e do CFR. O projeto mediano publica uma release a cada duas semanas, leva cerca de 11 dias do commit à release, tem CFR de CI de 4,7\% e se recupera de uma execução com falha em 2,6 horas. Na classificação geral, 87\% ficam em \textit{High} ou \textit{Medium} e só 8\% em \textit{Elite}.

A variante escolhida importa: o CFR (b) tem mediana quase quatro vezes maior que a do CFR (a), ao contrário do que esperávamos. A frequência de deploy tem correlação positiva com as duas variantes do CFR ($\rho$ de 0,16 e de 0,50), o que não confirma que velocidade e estabilidade andem juntas, embora o resultado com o CFR (b) seja em parte efeito da definição. Na RQ 06, só 4 dos 20 testes seguem significativos após a correção de Holm, e a idade é o fator mais associado ao desempenho. O \textit{rework rate} tem mediana de 31\% e correlação forte com o CFR (b) ($\rho$ de 0,71). Trocar as definições muda a categoria geral de 22\% a 30\% dos repositórios, quase sempre para uma vizinha: as conclusões gerais se mantêm, mas a categoria de cada repositório deve ser lida como estimativa.

O principal limite é a heurística de release corretiva, que sustenta o CFR (b) e o \textit{rework rate}: o seu F1 de 0,644 não atingiu a meta de 0,70, e as regras foram escolhidas sobre os mesmos itens da avaliação. Um próximo passo é validar a heurística em uma amostra maior e separada, e comparar as aproximações com dados reais de deploy e de incidentes.

\paragraph{Disponibilidade.} O \textit{pipeline}, os dados e as análises estão em \url{https://github.com/gabialvarenga/lab-experimentacao-software}, e o quadro de acompanhamento do projeto, em \url{https://github.com/users/gabialvarenga/projects/9}.
